\documentclass[12pt]{article}
\usepackage[a4paper, margin=1in]{geometry}
\usepackage{booktabs}
\usepackage{amsmath}

\title{COS 314 Assignment 2}
\author{
Jeandre Opperman 23542773 \\
Aidan Dawson 24593542
}
\date{}

\begin{document}

\maketitle

\section{Genetic Algorithm}

A Genetic Algorithm is a population based metaheuristic inspired by natural selection and evolution. It maintains a population of candidate solutions that evolve over generations using selection, crossover, and mutation. Fitter individuals are more likely to pass their traits, guiding the search toward optimal solutions.

\subsection{Special Relevance}

The 0/1 Knapsack Problem maps naturally to binary strings. GA performs strong global exploration, avoiding premature convergence and exploring multiple combinations simultaneously.

\subsection{Functionality}

\textbf{Representation and Initialization:}  
Solutions are binary strings where 1 means include and 0 means exclude. Initial population is randomly generated.

\textbf{Fitness and Constraint Handling:}  
Fitness is total value. If overweight, a stochastic repair removes items randomly until feasible.

\textbf{Selection:}  
Tournament selection with moderate pressure.

\textbf{Reproduction:}  
Single-point crossover with bit-flip mutation.

\textbf{Elitism:}  
Top individuals are preserved.

\subsection{Configuration}

\begin{itemize}
\item Population Size: 100
\item Generations: 500
\item Crossover Rate: 0.8
\item Mutation Rate: 0.02
\item Tournament Size: 5
\item Elitism Count: 2
\end{itemize}

\section{Iterated Local Search}

ILS is a trajectory-based metaheuristic that alternates between local search and perturbation to escape local optima.

\subsection{Special Relevance}

Knapsack landscapes are highly correlated. ILS exploits this with efficient local improvements and controlled exploration.

\subsection{Functionality}

\textbf{Initialization:}  
Random feasible solution.

\textbf{Local Search:}  
One-bit flip improvements until local optimum.

\textbf{Perturbation:}  
Three random bit flips.

\textbf{Acceptance:}  
Strict improvement only.

\subsection{Configuration}

\begin{itemize}
\item Iterations: 500
\item Neighborhood: Single bit flip
\item Perturbation Strength: 3
\end{itemize}

\section{Results}

\begin{center}
\begin{tabular}{lcccccc}
\toprule
Instance & Alg & Seed & Best & Optimum & Runtime \\
\midrule
f1 & ILS & 22 & 295 & 293 & 0.005 \\
   & GA  & 22 & 293 & 293 & 0.047 \\
f2 & ILS & 22 & 1024 & 1024 & 0.071 \\
   & GA  & 22 & 1024 & 1024 & 0.145 \\
f3 & ILS & 22 & 35 & 35 & 0.001 \\
   & GA  & 22 & 35 & 35 & 0.035 \\
f4 & ILS & 22 & 23 & 23 & 0.004 \\
   & GA  & 22 & 23 & 23 & 0.068 \\
f5 & ILS & 22 & 481.1 & 481.1 & 0.012 \\
   & GA  & 22 & 481.1 & 481.1 & 0.052 \\
f6 & ILS & 22 & 52 & 52 & 0.007 \\
   & GA  & 22 & 52 & 52 & 0.045 \\
f7 & ILS & 22 & 107 & 107 & 0.006 \\
   & GA  & 22 & 107 & 107 & 0.060 \\
f8 & ILS & 22 & 9767 & 9767 & 0.027 \\
   & GA  & 22 & 9767 & 9767 & 0.091 \\
f9 & ILS & 22 & 130 & 130 & 0.005 \\
   & GA  & 22 & 130 & 130 & 0.059 \\
f10 & ILS & 22 & 1025 & 1025 & 0.018 \\
    & GA  & 22 & 1025 & 1025 & 0.068 \\
knapPI & ILS & 22 & 9147 & 9147 & 0.203 \\
       & GA  & 22 & 7907 & 9147 & 0.341 \\
\bottomrule
\end{tabular}
\end{center}

\section{Statistical Analysis}

Wilcoxon signed-rank test ($\alpha = 0.05$).

Null hypothesis: equal performance.

Most instances (9/11) resulted in ties. Only 2 non-zero differences:

\begin{itemize}
\item f1: +2
\item knapPI: +1240
\end{itemize}

$W^+ = 3$, $W^- = 0$, $W = 0$.

Conclusion: insufficient sample size to reject null hypothesis.

\section{Analysis}

ILS consistently reaches high-quality solutions faster. The knapsack landscape favors greedy improvements, which ILS exploits efficiently. GA explores broadly but struggles with fine tuning and wastes time evaluating weak candidates early.

ILS proves more efficient and reliable for these instances.

\end{document}